\section{The Astronomy v1 Instance}
\label{sec:dataset}

Astronomy~v1 instantiates the protocol in exoplanet atmospheres---a
domain chosen because it uniquely combines a fast-moving literature,
structured catalogs, and space-telescope archives, and because the
2021--2026 window contains a natural experiment: JWST began delivering
data mid-window, resolving questions that were unanswerable at the
cutoff. Table~\ref{tab:corpus} summarizes the instance.

\begin{table}[t]
  \centering
  \small
  \caption{Astronomy v1 at a glance. Corpus manifests, frozen questions,
  retrieval records, and adjudicated labels are all released.}
  \label{tab:corpus}
  \begin{tabular}{@{}ll@{}}
    \toprule
    Historical cutoff & 2020-12-31 \\
    Domain scope & atmospheric composition in transmission spectra;
                   cloud/haze degeneracies; instrument systematics \\
    \midrule
    \corpusA{} (past) & 2{,}512 deduplicated papers, 2015--2020 (NASA ADS;
                        500-paper full-text core) \\
    \corpusB{} (future) & 1{,}891 unique papers, 2021--2026, temporally
                          isolated \\
    Questions & 10, frozen, ranked, with pre-cutoff source evidence \\
    Retrieval & \texttt{text-embedding-3-small}, cosine, top-8; records released \\
    Judge & \texttt{gpt-4.1}, temperature 0, citation-constrained;
            human-adjudicated \\
    \bottomrule
  \end{tabular}
\end{table}

\paragraph{Corpora.} Both corpora are defined by frozen manifests---ADS
query sets, date windows, deduplication and filtering rules
(records without abstracts are dropped)---rather than by bulk data
dumps: the manifests are committed, and a script rebuilds either corpus
from its manifest via the ADS API. \corpusB{}'s manifest adds targeted
follow-up queries for the questions' objects (HD~189733, HD~209458,
WASP-12, WASP-121, TRAPPIST-1) so that engagement is measured against
the relevant future literature rather than against whatever a generic
query happens to return.

\paragraph{Leakage controls.} Temporal isolation is enforced
mechanically, not editorially: (1) nothing dated after the cutoff may
enter \corpusA{}, including catalog rows updated post-cutoff; (2) every
question's source evidence must predate the cutoff; (3) every retrieved
document must postdate it; (4) the judge may cite only retrieved
candidates; (5) question/retrieval/annotation records must align
one-to-one; (6) \corpusB{}'s window must start strictly after the
cutoff. All six checks run in continuous integration on every change to
the released data, together with a check that the released
\texttt{metrics.json} reproduces bit-identically from the raw
annotations.

\paragraph{Questions.} The ten frozen questions were generated by an
evidence-graph system \citep{paper1} from \corpusA{} only: claims with
provenance are extracted from the full-text core, cross-paper tensions
are detected and typed (observational tensions, methodological
challenges, single-dataset conclusions, independent qualifications), and
surviving signals are refined into ranked, falsifiable questions. For
the benchmark, that system is submission
\texttt{evidence\_graph\_v1}---evaluated through the same interface as
any future submission. Each released record carries the question text,
system rank, signal type, target objects, and pre-cutoff source
bibcodes; a curation log documenting human edits made \emph{before}
freezing (including a near-duplicate merge and presupposition fixes,
Section~\ref{sec:errors}) is released for audit.

\paragraph{What is released.} Frozen questions; corpus manifests; per-
question retrieval records (model, window, top-$k$, judge-cited
documents, top-1 similarity; full ranked lists with scores are scheduled
for v1.1); adjudicated outcome annotations with rationales and
supporting bibcodes; the adjudication log; computed metrics and
per-question results; four frozen baseline submissions with their full
retrieval records (per-document scores included), judge annotations,
and reports (Section~\ref{sec:baselines}); the scaled v1L instance
(manifests, configs, 424 frozen baseline questions, retrieval records
with per-document scores, judge annotations, per-system reports, and
the statistical comparison script of Section~\ref{sec:scaled}); the
tension-pair generators, the four-cutoff stress-test corpora manifests,
all 798 stress-test judgments, and the deterministic specificity rubric
of Section~\ref{sec:foresight}; and the complete pipeline code with
tests. The
data format is deliberately plain (JSONL + JSON manifests) so that other
groups can mint new instances---different domain, different cutoff---by
writing two manifests and one config file.
